\documentclass[11pt]{article}
\usepackage{fontspec}
\usepackage[french]{babel}
\usepackage{amsmath}
\usepackage{amssymb}
\usepackage{amsthm}
\usepackage{mathrsfs}
\usepackage{tikz-cd}
\usepackage{enumitem}
\usetikzlibrary{babel}

\newcommand{\printedpage}[1]{\marginpar{\footnotesize p.~#1}}
\newcommand{\msp}[1]{\marginpar{\footnotesize [#1]}}
\newcommand{\X}{\mathcal{X}}

\begin{document}

\printedpage{14}
% head: A. Grothendieck : Dérivateurs, Ch. II, Édition des Universités de Montpellier II et Paris VII
% foot: Version 8 juillet 2001

\[
\begin{tikzcd}
\mathrm{b}' \arrow[r, Rightarrow] & \mathrm{b} \arrow[rr, Rightarrow] \arrow[d, Rightarrow] & & \mathrm{a} \arrow[d, Rightarrow] \\
 & \mathrm{d} \arrow[r, Leftrightarrow] & \mathrm{e} \arrow[r, Leftrightarrow] & \mathrm{c}.
\end{tikzcd}
\]

Prouvons les implications de la ligne du bas. Considérons d). Si $\begin{tikzcd}[cramped, row sep=0.1em, column sep=0.8em, baseline=(current bounding box.center)] a \arrow[dr] & \\ & x \\ b \arrow[ur] & \end{tikzcd}$ \texttt{[est]} dans $S$, dire que le produit fibré dans $\X$ est `non vide' signifie, puisque $S$ est un système de générateurs et en particulier \textit{non vide,} qu'il existe $c$ dans $S$ et $c\longrightarrow a\times_x b$, ou ce qui revient au même, $\begin{tikzcd}[cramped, row sep=0.1em, column sep=0.8em, baseline=(current bounding box.center)] & a \\ c \arrow[ur, dotted] \arrow[dr, dotted] & \\ & b \end{tikzcd}$, rendant commutatif le diagramme $\begin{tikzcd}[cramped, row sep=0.1em, column sep=0.8em, baseline=(current bounding box.center)] & a \arrow[dr] & \\ c \arrow[ur, dotted] \arrow[dr, dotted] & & x \\ & b \arrow[ur] & \end{tikzcd}$. Donc d) équivaut à PS 1, i.e. à e) ($\X$ localement faiblement filtrant à gauche). Considérons la condition c) : les $\X/a=(S/a)^\sim$ sont irréductibles. Cela signifie que pour deux sous-objets \textit{non vides} $U,V\subseteq a$ dans $\X^\wedge$, on a $U\cap V\neq\emptyset$, ou encore que pour deux cribles

\msp{page 22}$\underline{U}$, $\underline{V}$ de $S/x$, si $\underline{U}$, $\underline{V}$ non vides (ce qui signifie que les sous-objets de $x$ dans $\X^\wedge$ qu'ils définissent sont `non vides', ou que les objets de $S$ sont `non vides' dans $\X^\wedge$), alors $\underline{U}\cap\underline{V}$ est non vide. On peut supposer que $\underline{U}$, $\underline{V}$ est engendré chacun par une seule flèche $a\xrightarrow{u}x$, $b\xrightarrow{v}x$. Le crible intersection est formé des objets $c\longrightarrow x$ de $S/x$ tels que $c\longrightarrow x$ se factorise par un $\begin{tikzcd}[cramped, row sep=0.1em, column sep=0.8em, baseline=(current bounding box.center)] & a \\ c \arrow[ur, dotted] \arrow[dr, dotted] & \\ & b \end{tikzcd}$, donc dire que ce crible n'est pas vide, signifie que PS 1 est satisfait pour les deux flèches données $u$, $v$. Donc c) équivaut à la condition PS 1\texttt{[, i.e. à]} e.

Ainsi on a établi
\[
\mathrm{b}'\Longrightarrow\mathrm{b}\Longrightarrow\mathrm{a}\Longrightarrow(\mathrm{c}\Longleftrightarrow\mathrm{d}\Longleftrightarrow\mathrm{e}),
\]
reste à prouver que c $\Longrightarrow$ b$'$. Soit donc $\begin{tikzcd}[cramped, row sep=0.1em, column sep=0.8em, baseline=(current bounding box.center)] U \arrow[dr] & & V \arrow[dl] \\ & W & \end{tikzcd}$ un diagramme d'objets `non vides', avec $W$ 0-connexe, il faut prouver que le produit fibré est non vide. Considérant $a\longrightarrow U$ et $b\longrightarrow V$ avec $a,b\in S$, on est ramené au cas où $U=a$, $V=b$ dans $S$. Prouvons que si $a\times_W b$ est `vide', alors $W$ est disconnexe. Pour ceci, considérons la sous-catégorie $C$ de $S/W$ formée des $x\longrightarrow W$ tels que $a\times_W x$ soit `non vide'. Désignant par $w(a)$, $w(x)$ l'image de $a$, $x$ dans $W$, cela signifie que $w(a)\cap w(x)\neq\emptyset$.

\msp{page 23}Notons que tout quotient dans $\X^\wedge$ d'un objet irréductible est irréductible, donc $w(a)$, $w(x)$ sont des objets irréductibles. Cela montre que pour $x$ dans $C$ et $x'$ dans $S/W$, on a $x'$ dans $C$, i.e. $w(a)\cap w(x')\neq\emptyset$ si et seulement si $w(x)\cap w(x')\neq\emptyset$. (Lemme direct : dans un objet $W$ d'un topos $\X$, si on a \texttt{[des]} sous-objets $W_1$, $W_2$, $W_3$ irréductibles, si l'un des deux \texttt{[plutôt trois]} rencontre les deux autres, ceux-ci se rencontrent.) Cela montre que $C$ est un crible, puisque $x'\longrightarrow x$ implique $w(x')\subseteq w(x)$, donc $w(x')$ rencontre $w(x)$, donc aussi $w(a)$. Mais alors le co-crible $C'=(S/W)\backslash C$ est également un crible, car si $c$
\printedpage{15}
% head: A. Grothendieck : Dérivateurs, Ch. II, Édition des Universités de Montpellier II et Paris VII
% foot: Version 8 juillet 2001
est dans $C'$, i.e. \texttt{[si]} $w(c)$ ne rencontre pas $w(x)$, a fortiori pas tout $x'$ au-dessus de $x$, pour lequel on aura $w(x')\subseteq w(x)$. Ainsi, $S/W$ est décomposé en somme de deux ouverts $C$, $C'$, non vides puisque l'un contient $a$, l'autre $b$. Par suite, $W$ est décomposé en somme de deux sous-objets non vides, donc $W$ est non 0-connexe.

On a prouvé 2.8, donc aussi (revenant à 2.6) les implications vers PS 1 dans le diagramme page 20. Il reste à prouver
\[
\left\{
\begin{aligned}
&\mathrm{b}\Longrightarrow\mathrm{PS}\ 2,\ \mathrm{c}'\Longrightarrow\mathrm{PS}\ 2\\
&\mathrm{PS}\ 1+\mathrm{PS}\ 2\ (\text{i.e. d})\Longrightarrow\mathrm{a}'
\end{aligned}
\right.
\]
(\footnote[32]{NB finalement on ne prouve pas c$'\Longrightarrow$ PS 2.}).

\msp{page 24}Nous nous plaçons donc dans l'hypothèse où PS 1 est satisfait (puisque b), c$'$) et d) impliquent PS 1), donc où on est sous les conditions équivalentes de 2.8 - $\X$ est localement irréductible. Considérons alors la condition PS 2, et prouvons qu'elle est impliquée tant par b), que par c$'$), et qu'elle implique a$'$). Elle signifie que pour une double flèche $a\overset{u}{\underset{v}{\rightrightarrows}}b$ dans $S$, le conoyau \texttt{[plutôt noyau]} dans $S^\wedge$, ou, ce qui revient au même, dans $\X^\wedge=S^\sim\subseteq S^\wedge$, est `non vide'. On peut l'écrire sous la forme d'un produit fibré
\[
\begin{tikzcd}
a \arrow[r, "{\scriptstyle (u,v)}"] & b\times b \\
\operatorname{Ker}(u,v) \arrow[u] \arrow[r] & b. \arrow[u, "\mathrm{diag}"']
\end{tikzcd}
\]
Ici par hypothèse $a$ et $b$ sont des objets 0-connexes de $\X^\wedge$ (ceux de $S$ l'étant), c) est vérifié (($\X$ localement irréductible et) ses composantes connexes \texttt{[sont]} totalement 0-connexes), alors $b\times b$ est 0-connexe comme produit de deux objets 0-connexes sur une \textit{même composante connexe,} laquelle est totalement 0-connexe. Comme $\X\times\X$ est localement irréductible, il s'ensuit que le produit fibré est non vide (2.8), donc b) $\Longrightarrow$ d). Prouvons d) $\Longrightarrow$ a), i.e. si

\msp{page 25}$S$ est localement cofiltrante, alors les topos induits $\X/W$, où $W$ est 0-connexe, sont totalement 0-connexes. Ce topos induit est décrit par le site $S/W$, i.e. $S/W$ est un système de générateurs. Or si $S$ est localement filtrante, $S/W$ (étale sur $S$) l'est aussi. Donc on peut supposer $W=e$, et à prouver que si $S$ est localement cofiltrante et $\X$ 0-connexe, alors $\X$ est totalement 0-connexe. Notons que $\X$ 0-connexe équivaut à $S$ 0-connexe - il faudrait quand même l'expliciter :

\textbf{Lemme 2.9.} \textit{Soit $\X$ un topos, $S\subseteq\X^\wedge$ une sous-catégorie pleine génératrice formée d'objets} 0-connexes. \textit{Alors $\X$ est $0$-connexe si et seulement si $S$ l'est - plus généralement, l'ensemble ordonné des sous-objets à la fois ouverts et fermés de $S$ dans} Cat \textit{(ou encore des
\printedpage{16}
% head: A. Grothendieck : Dérivateurs, Ch. II, Édition des Universités de Montpellier II et Paris VII
% foot: Version 8 juillet 2001
sous-objets ouverts $=$ cribles de $S$ qui sont aussi fermés, i.e. des cocribles) est isomorphe à celui des sommandes directs de $\X$.}

Je l'ai utilisé tacitement dans la démonstration de 2.8, pour le topos induit $\X/W$. Je n'explicite pas la démonstration.

Ainsi (revenant à la démonstration de d) $\Longrightarrow$ c) dans 2.6.) $S$ étant 0-connexe et localement cofiltrante,

\msp{page 26}\texttt{[elle]} est en fait cofiltrante. Donc pour deux objets $a$, $b$ de $S$, $\exists\ c$ avec $c\longrightarrow a$, $c\longrightarrow b$ (condition PS 1$'$, plus forte que PS 1). Pour prouver que $\X$ est totalement 0-connexe, on est réduit à prouver par 2.3 que les $a\times b$ dans $\X^\wedge$ sont 0-connexes. Par 2.9, il revient au même de prouver qu'ils sont 0-connexes dans $S^\wedge$. Donc on est ramené à prouver le

\textbf{Lemme 2.10.} \textit{Soit $S$ une catégorie cofiltrante. Si $a$, $b$ sont deux objets de $S$, $a\times b$ est un objet} $0$\textit{-connexe de $S^\wedge$.}

On voit de suite qu'il suffit de prouver (et en fait il revient au même, compte tenu que $a\times b$ est localement irréductible) que si on a deux flèches
\[
\begin{tikzcd}
x \arrow[r] & a\times b \\
 & y, \arrow[u]
\end{tikzcd}
\]
alors $x\times_{a\times b}y\neq\emptyset$, i.e. $\exists\ z$ dans $S$ rendant commutatif
\[
\begin{tikzcd}
x \arrow[r] & a\times b \\
z \arrow[u] \arrow[r] & y. \arrow[u]
\end{tikzcd}
\]
Or par PS 1$'$ il existe bien un $z$ s'envoyant dans $x$ et dans $y$, mais les deux flèches composées
\[
z\rightrightarrows a\times b,
\]
i.e. chacune des double flèches
\[
z\rightrightarrows a \tag*{$(*)$}
\]
\[
z\rightrightarrows b \tag*{$(*')$}
\]
peuvent n'être pas égales. Mais par PS 2,

\msp{page 27}on peut trouver
\[
z'\longrightarrow z
\]
\printedpage{17}
% head: A. Grothendieck : Dérivateurs, Ch. II, Édition des Universités de Montpellier II et Paris VII
% foot: Version 8 juillet 2001
égalisant la double flèche $(*)$, donc on peut supposer que celle-ci est déjà triviale. Appliquant encore une fois PS 2, par un $z''\longrightarrow z'$, on égalise de même la double flèche $(*')$, et on gagne. On a donc prouvé, finalement,
\[
\begin{tikzcd}
\mathrm{a} \arrow[r, Leftrightarrow] \arrow[d, Rightarrow] & \mathrm{a}' \arrow[r, Rightarrow] & \mathrm{b} \arrow[r, Leftrightarrow] & \mathrm{d}\ (=\mathrm{PS}\ 1+\mathrm{PS}\ 2) \\
\mathrm{c} \arrow[r, Leftrightarrow] & \mathrm{c}',
\end{tikzcd}
\]
mais pas l'implication inverse. En d'autres termes, si les topos induits $\X/a$, pour $a$ dans $S$, sont totalement 0-connexes, je n'ai pas prouvé qu'il en est de même des $\X/W$, $W$ un objet 0-connexe \textit{quelconque} de $\X^\wedge$. Ainsi, on en est au point suivant

\textbf{Perplexité 2.11.} \textit{Soit $X$ dans} Cat. \textit{Si $X$ est localement filtrante à gauche, alors les produits fibrés $a\times_x b$ d'objets dans $X$ sont des objets} $0$\textit{-connexes de $X^\wedge$, et cette condition à son tour implique PS 1, i.e. que $X$ est localement irréductible. Mais en présence de cette hypothèse PS 1, ou $X$ localement irréductible, la condition sur les produits fibrés $a\times_x b$ équivaut à la suivante : pour tout diagramme} \texttt{[en]} \textit{traits pleins}

\msp{page 28}
\[
\begin{tikzcd}
c \arrow[r] & a\times_x b \\
e \arrow[u, dotted] \arrow[r, dotted] & d, \arrow[u]
\end{tikzcd}
\qquad\qquad
\textit{avec $x$, $a$, $b$, $c$, $d$ dans $X$,}
\]
\textit{on peut compléter en un diagramme} commutatif, \textit{avec $e$ dans $X$. En d'autre termes, tout diagramme commutatif à traits pleins ci-dessous se complète en y rajoutant un `objet initial' $e$} % [sic] d'autre termes -> d'autres termes, missing plural s
\[
\begin{tikzcd}[column sep=small, row sep=small]
 & & a \arrow[ddr] & \\
 & c \arrow[ur] \arrow[dddr] & & \\
e \arrow[ur, dotted] \arrow[dr, dotted] & & & x \\
 & d \arrow[uuur] \arrow[dr] & & \\
 & & b \arrow[uur] &
\end{tikzcd}
\]
(\footnote[33]{Il suffit d'exiger pour $e$ la commutativité des deux carrés $(e\ c\ b\ d)$ et $(e\ d\ a\ c)$.}).

Cette condition est donc impliquée par $X$ localement cofiltrante, et elle est satisfaite aussi par ce qui précède plus généralement si et seulement si (\footnote[34]{si et seulement si ?}) les $X/x$ sont cofiltrantes (\footnote[35]{Elle signifie que $X$ satisfait PS 1, et la condition

\textbf{PS$'$ 2} Tout diagramme commutatif
\[
a\rightrightarrows b\longrightarrow x
\]
% footnote continued on printed page 18
s'insère dans un diagramme commutatif
\[
y\longrightarrow a\rightrightarrows b\longrightarrow x.
\]
}). Mais la réciproque est-elle vraie - cette condition de totale 0-connexité des $X/x$ implique-t-elle que $X$ soit totalement 0-connexe ? Mais je n'ai pas construit d'exemple où les $X/x$
\printedpage{18}
% head: A. Grothendieck : Dérivateurs, Ch. II, Édition des Universités de Montpellier II et Paris VII
% foot: Version 8 juillet 2001
sont cofiltrantes sans que $X$ soit localement cofiltrante (ou ce qui revient au même, de $X$ 0-connexe avec les $X/x$ cofiltrantes sans que $X$ en soit cofiltrante)(\footnote[36]{Mais il y a un \textit{contre-exemple} évident, si $X$ est un groupoïde 0-connexe, les $X/x$ sont $\approx e$, donc cofiltrantes, mais $X$ n'est cofiltrant (i.e. \texttt{[mot(s) illisible(s)]}) que si $X$ est 1-connexe. P.ex. si $C$ est un groupoïde $\neq 1$, alors $X=B_G$ est tel que les $X/x$ sont cofiltrantes, mais $X$ non cofiltrant. % [sic] si $C$ est un groupoïde -> si $G$ est un groupoïde, the groupoid is then written $B_G$
}) ; ni d'exemple où les $X/x$ sont totalement 0-connexes, sans que les $X/x$ soient cofiltrantes. Mais notons le

\textbf{Corollaire 2.12 :} \textit{Les conditions équivalentes a), a$'$), b), d) (à l'exclusion de c), c$'$), plus faibles apparemment) du théorème 2.6}

\msp{page 29}\textit{équivalent encore à la suivante}
\begin{itemize}[leftmargin=2.5em]
\item[\textit{e)}] \textit{$\forall\ x$ dans $S$, le topos induit $\X/x$ et le topos induit $\X/x\times x$ est irréductible} (\footnote[37]{NB $\X/x$ inutile.}),
\end{itemize}
\textit{où encore à la suivante} % [sic] où encore -> ou encore, conjunction "ou" without accent
\begin{itemize}[leftmargin=2.5em]
\item[\textit{e$'$)}] \textit{$\forall\ x$ dans $S$, $S^\wedge/x$ et $S^\wedge/x{\times}x$ sont irréductibles, i.e. $S/x$ et $S/x{\times}x$ sont irréductibles} (\footnote[38]{$S/x$ inutile, car $x\xrightarrow{\mathrm{diag}}x\times x$ est une immersion ouverte non vide, donc si $x\times x$ est irréductible, $x$ l'est aussi. Mais il revient au même de dire que les $S/x$ sont irréductibles (condition PS$^*$ 1, cf. 2.8, p. 21) \textit{et} que les $x\xrightarrow{\mathrm{diag}}x\times x$ sont dominants (ce que implique en effet que $x\times x$ est irréductible). Ceci redonne immédiatement la forme PS$^*$ 2 de la condition cofiltrante. % [sic] ce que implique -> ce qui implique, relative pronoun subject
}).
\end{itemize}
C'est la forme `topologique' la plus faible en apparence que je peux trouver pour la condition $S$ localement cofiltrante, i.e. $\X$ (ou Top$(S)$, au choix) localement totalement 0-connexe.

\textbf{Définition 2.13.} On dit qu'un topos $\X$ est \textit{localement totalement} $0$\textit{-connexe} s'il satisfait aux conditions équivalentes a, a$'$, b, d du théorème 2.6, i.e. si $\X$ est localement connexe et si de plus tout topos induit 0-connexe est totalement 0-connexe - le produit fibré de deux objets 0-connexes sur un troisième est 0-connexe. Il revient au même de dire que toute (ou que une) sous-catégorie pleine génératrice $S$ de $\X^\wedge$ formée d'objets 0-connexes est localement cofiltrante (en supposant toujours $\X$ localement connexe, ou ce qui revient au même, l'existence d'un $S$ formé d'objets 0-connexes).

On n'a pas terminé encore la démonstration de 2.6 - il reste à prouver que si $\X$ (supposé localement totalement 0-connexe) est 0-connexe, il est $\infty$-connexe - ou encore que les composantes connexes de $\X$, supposé localement totalement 0-connexe, sont $\infty$-connexes, i.e. $W_\infty$-asphériques.

\msp{page 30}\textbf{Corollaire 2.14.} \textit{Notons quand même que si $\X$ est de la forme Top$(T)$, $T$ un espace topologique, alors $\X$ est localement totalement} $0$\textit{-connexe si et seulement si il est localement irréductible, i.e. si et seulement si $T$ est localement irréductible.}

\end{document}
